% ------------------------------------------------------------------------------
% Package Dasar
% ------------------------------------------------------------------------------
\usepackage[utf8]{inputenc}
\usepackage{sectsty}
\usepackage{graphicx}
\usepackage{epstopdf}
\usepackage{algorithm}
\usepackage{algpseudocode}
\usepackage{array}
\usepackage[table]{xcolor}
\usepackage{amsmath}
\usepackage{amssymb}
\usepackage[bahasa]{babel}
\usepackage{indentfirst}

% ------------------------------------------------------------------------------
% Tata Letak Halaman
% ------------------------------------------------------------------------------
% Margin: kiri 4cm, kanan-atas-bawah 3cm
\usepackage[top=3cm, bottom=3cm, left=4cm, right=3cm]{geometry}

\usepackage{multirow}
\usepackage{hhline}

% ------------------------------------------------------------------------------
% Sitasi dan Referensi
% ------------------------------------------------------------------------------
\usepackage{etoolbox}
\usepackage[numbers]{natbib}
\usepackage[bookmarks,hypertexnames=false,debug]{hyperref}
\usepackage{bookmark}

% ------------------------------------------------------------------------------
% Penyesuaian Format Bab dan Diagram
% ------------------------------------------------------------------------------
\makeatletter
\patchcmd{\@makechapterhead}{\vspace*{50\p@}}{\vspace*{0pt}}{}{}
\patchcmd{\@makeschapterhead}{\vspace*{50\p@}}{\vspace*{0pt}}{}{}
\makeatother

\chapterfont{\centering}
\newcommand{\bigsize}{\fontsize{16pt}{14pt}\selectfont}
\chapterfont{\centering\bigsize\bfseries}
\sectionfont{\large\bfseries}
\usepackage{tikz}
\usetikzlibrary{shapes.geometric, arrows}

% ------------------------------------------------------------------------------
% Penomoran: Bab, Subbab, Persamaan, Gambar, Tabel
% ------------------------------------------------------------------------------
\renewcommand{\thechapter}{\Roman{chapter}}
\renewcommand\thesection{\arabic{chapter}.\arabic{section}}
\renewcommand\thesubsection{\thesection.\arabic{subsection}}
\renewcommand{\theequation}{\arabic{chapter}.\arabic{equation}}
\renewcommand{\thefigure}{\arabic{chapter}.\arabic{figure}}
\renewcommand{\thetable}{\arabic{chapter}.\arabic{table}}

% ------------------------------------------------------------------------------
% Header dan Footer
% ------------------------------------------------------------------------------
\renewcommand\bibname{Daftar Pustaka}
\addto{\captionsbahasa}{
  \renewcommand{\bibname}{DAFTAR PUSTAKA}
  \renewcommand{\listfigurename}{DAFTAR GAMBAR}
  \renewcommand{\listtablename}{DAFTAR TABEL}
  \renewcommand{\contentsname}{DAFTAR ISI}
  \renewcommand{\chaptername}{BAB}
}
\usepackage{fancyhdr}
\pagestyle{fancy}
\lhead{}
\chead{}
\rhead{}
\lfoot{}
\cfoot{\thepage}
\rfoot{}
\renewcommand{\headrulewidth}{0pt}

% ------------------------------------------------------------------------------
% Kustomisasi Daftar Isi
% ------------------------------------------------------------------------------
\usepackage[titles]{tocloft}
\renewcommand{\cftchappresnum}{BAB }
\settowidth{\cftchapnumwidth}{BAB VIII }

% ------------------------------------------------------------------------------
% Konfigurasi Daftar Lampiran Otomatis
% ------------------------------------------------------------------------------
\makeatletter
\newcommand\listofappendices{%
  \chapter*{DAFTAR LAMPIRAN}%
  \@starttoc{loa}%
}
\newcommand\lampiran[1]{%
  \chapter*{#1}%
  \addcontentsline{loa}{chapter}{#1}%
}
\makeatother
